\begin{document}
\docTitle{Custom commands in manual\_style}
\phantomsection
\label{doc:manual-style}
\ReviewAnchor{mjpdf-922d6876-a70a-4ed2-931e-7f0c60c8aa8a}{We use a combination of standard and custom commands...}{This documents describes custom shared LaTex commands available for the note syst}ems. 
\agentProposedEdit{agent-remove-6b31f515-6699-4f66-81bb-8ac1de88d391}{each project now owns the style used by its active TeX dependency graph}{These commands are implemented through the file \path{tools/latex/manual_style.sty}, the shared package loaded by the standard renderer.}{}
\agentProposedEdit{agent-remove-a792028b-e21b-4f25-a910-450a773f851b}{this document now describes the project-local style instead of the retired workspace-wide path}{This document is the public command reference for \path{tools/latex/manual_style.sty}.}{}
\agentProposedEdit
  {agent-add-fb7f7740-c63c-5b6f-a5f9-bfc122e0ffd5}
  {Add this agent-authored block for review.}
  {}
  {For \code{docs-template}, these commands are implemented by the project-owned package at \path{tex/support/manual_style.sty}.
Every active \code{docs-*} project must load its own required style from \path{tex/support/} so the exact build input can be captured, hashed, published, and reconstructed independently.}

\phantomsection
\section*{Linking}\label{sec:manual-style-tex-link}
\subsection*{Within Tex documents}
Use the standard \code{.tex} command

\begin{DocCode}
\hyperref[sec:parallel-execution]{the Parallel Execution section}
\end{DocCode}

Define the target as:

\begin{DocCode}
\phantomsection
\subsection*{Parallel Execution}
\label{sec:parallel-execution}
\end{DocCode}

\subsection*{To a different .tex documents}
Use \code{\TexLink}.
Creates a link to a labelled location in another renderable source document.
Its arguments are the logical target \code{docs-project/file.tex}, the target's LaTeX label, and the visible link text.

\begin{DocCode}
% In docs-concepts/tex/source/sdpo.tex:
\phantomsection
\subsection*{Implementation details}
\label{sec:sdpo-implementation}

% In another source document:
\TexLink
  {docs-concepts/sdpo.tex}
  {sec:sdpo-implementation}
  {the SDPO implementation section}
\end{DocCode}

The selector is intentionally independent of the generated PDF filename.
The target must be a standalone source under the selected project's \path{tex/source/} folder, and its label must exist in the rendered \code{.aux} file.
Use \code{\phantomsection} before a labelled unnumbered heading so the link resolves to the heading's location.
The standard renderer turns the link into an exact external jump in an individual PDF.
For a combined PDF, it appends a cross-project target once and converts the link into an internal jump.
Use \code{\href} for ordinary web URLs and \code{\TexLink} for links between documents.
Do not use the \code{run:} prefix: it creates an application-launch action that many PDF viewers block.

\subsection*{To a PDF document with no .tex}
Use \code{\PdfLink}.
The PDF must be located at PDF stored directly in the workspace-global \path{third_party/} folder.
The paper key is the PDF filename without its \code{.pdf} suffix, and the anchor name identifies one annotation whose complete comment is \code{doc-anchor:anchorName}.
\agentProposedEdit
  {agent-add-f982447d-06af-5a67-b1d6-5305bab41bed}
  {Add this agent-authored block for review.}
  {}
  {The PDF filename must start with a letter or digit and otherwise contain only letters, digits, dots, underscores, or hyphens.}
\agentProposedEdit{agent-remove-e4186ade-9a4b-464c-95c0-8d58f1a988c1}{replace the former Android application name}{The MJ PDF \textbf{Put in Library} workflow enforces this filename rule and preserves the imported PDF in the workspace-global \path{third_party/} folder.}{}
\agentProposedEdit
  {agent-add-3900fb10-2e1c-5a8e-977c-9ffebf830230}
  {Add this agent-authored block for review.}
  {}
  {The Thoughts Android \textbf{Put in Library} workflow enforces this filename rule and preserves the imported PDF in the workspace-global \path{third_party/} folder.}
The renderer jumps to the annotation geometry in an individual PDF and preserves the exact position after appending the paper to a combined PDF.
\agentProposedEdit
  {agent-add-29471f31-0fcc-5296-af27-4f7e711e24eb}
  {Add this agent-authored block for review.}
  {}
  {A combined PDF appends each referenced \code{\PdfLink} target and each cross-project \code{\TexLink} target once, while same-project \code{\TexLink} targets are already present in manifest order.}
Use \code{\TexLink} with a LaTeX label whenever the target has renderable workspace source.


\ReviewAnchor{mjpdf-256f2140-6cbd-4f02-8df7-760c7167cadb}{Add quotes?}%
\section*{Document title and source name}

\paragraph{\code{\docTitle}.}
Sets the document name used in the page header, configures the header and footer, enables configured line numbering, and normally prints a centered title.
An optional first argument supplies a shorter header name.

\begin{DocCode}
\docTitle{Full document title}
\docTitle[Short header]{Full document title}
\end{DocCode}

\paragraph{\code{\code}.}
Prints literal inline text in a monospaced font without interpreting LaTeX commands inside the argument.

\paragraph{\code{DocCode}.}
A verbatim block for commands, paths, logs, and source examples.
Long lines may wrap.
When document line numbers are enabled, its numbering continues from the surrounding document.
Commands written inside this environment are displayed, not executed.
\agentProposedEdit
  {agent-add-5179bacf-5ee5-5414-b275-29929868ad28}
  {Add this agent-authored block for review.}
  {}
  {Use \code{\begin{DocCode}[numbered]} to display independent one-based line numbers for one block in every PDF while leaving plain blocks unnumbered.}

\paragraph{\code{LinkedDocCode}.}
A verbatim-style block that keeps long lines unwrapped and enables backslash and braces as command characters.
This allows selected pieces of a displayed line to contain active LaTeX formatting or links.
It shares the continuous line-number behavior of \code{DocCode}.
\agentProposedEdit
  {agent-add-d36f3ba3-95bb-5ac1-b95c-6f83e0ad36f8}
  {Add this agent-authored block for review.}
  {}
  {The same opt-in syntax is available as \code{\begin{LinkedDocCode}[numbered]}.}

\begin{LinkedDocCode}
docs-concepts/tex/source/
|-- \TexLink{docs-concepts/sdpo.tex}{sec:sdpo-implementation}{sdpo.tex}  \textit{(linked target)}
`-- README.md
\end{LinkedDocCode}

\paragraph{\code{DocComment}.}
Formats a short document comment as a small gray italic quotation prefixed with ``Comment.''

\section*{Source-level review anchors}
\label{sec:manual-style-review-anchor}

\paragraph{\code{\ReviewAnchor}.}
Defines the source-level marker used by the SyncTeX comment-anchoring workflow.
\agentProposedEdit{agent-remove-1903a283-fb52-463c-bb93-3091570ff1e2}{the active command is provided by each project's exact project-owned style input}{The command is implemented in \path{tools/latex/manual_style.sty} and is available to every standard \code{docs-*} source rendered through the shared package.}{}
\agentProposedEdit
  {agent-add-8ee596b4-e22e-500e-b839-e6136fd26c7a}
  {Add this agent-authored block for review.}
  {}
  {The command is implemented in this project's \path{tex/support/manual_style.sty}, and other standard projects provide the same interface from their own project-owned support package.}

The command takes a stable comment identifier followed by a readable mirror of
the actual comment:
\begin{DocCode}
\ReviewAnchor{example-comment-id}{The system should regenerate all PDFs.}%
The system periodically regenerates all PDFs.
\end{DocCode}

The command must produce no visible text and must not alter spacing or line breaking.
During compilation it emits a uniquely named PDF destination at its source position.
The comment extraction and retention tools use that destination to recreate the associated PDF annotation after text reflows, paragraphs move, or earlier pages are added or removed.

The second argument is a one-line, LaTeX-safe editable comment for people
reading the source; it produces no output.
Whitespace is collapsed, but the text is not truncated.
\agentProposedEdit{agent-remove-30db119b-20a3-4b96-a7cb-c5129878f413}{replace the former desktop application name}{Changing it updates the canonical \path{comments/*.jsonl} record and the open
Okular annotation; Okular and companion edits update it in the other
direction.  The JSONL record remains authoritative for author, timestamps,
annotation style, selected text, source context, and synchronization state.}{}
\agentProposedEdit
  {agent-add-7a9f62e5-444b-5514-aef5-84605ced5862}
  {Add this agent-authored block for review.}
  {}
  {Changing it updates the open Thoughts Desktop annotation and the local materialized \path{comments/*.jsonl} record, then queues the corresponding durable operation for the next server synchronization.
PostgreSQL remains authoritative for mutable comment and review projections, while JSONL records are materialized interoperability, recovery, anchoring, and PDF-publication views.}
Per-comment common-ancestor metadata resolves simultaneous text edits by
timestamp, with ties favoring the canonical record.

\agentProposedEdit{agent-remove-73d51495-5ce2-494a-9f03-b5a789d9a426}{replace the former desktop application name}{\code{\ReviewAnchor} markers are inserted by Okular Enter's shared native synchroniser, either during an interactive save or through its headless pre-render helper, rather than written manually.}{}
\agentProposedEdit
  {agent-add-2cb85006-529e-5827-9545-49167c170c5a}
  {Add this agent-authored block for review.}
  {}
  {\code{\ReviewAnchor} markers are inserted by the Thoughts Desktop durable mutation worker, the in-process source-anchor watcher, or the shared headless pre-render helper rather than written manually.
Thoughts Android mutations are synchronized through Fedora and are materialized on the Mac before the same native anchoring rules run.}
The destination may be any authored \code{.tex} input inside the same \code{docs-*} project when that file is present in the exact build's \code{.fls} manifest; it is not restricted to the top-level \path{tex/source/} wrapper.
The first document-level \code{\docTitle} after
\code{\textbackslash begin\{document\}} is a hard placement floor.
If reverse SyncTeX maps a comment before that title or onto any line occupied
by the complete title command, including an optional or multiline argument,
the marker is inserted immediately after the title's closing line.
No review marker is inserted above \code{\docTitle}.
Other unsafe structural commands remain \code{needs_anchor}; the synchroniser
does not search for an arbitrary nearby line.
Markers are not inserted in generated, build-output, PDF, archive, third-party, or version-control directories.
Each identifier must occur at most once in the managed source tree.
Repeated comment extraction must reuse the existing marker instead of creating a duplicate.
Moving or renaming a source file is safe as long as the unique marker moves with its associated passage.
During repository-wide preparation, an unanchored comment can also follow a
missing source into another standard \code{docs-*} project when its exact-build
snapshot supplies a normalized target-and-context signature that occurs in
exactly one current managed source.
The comment keeps its stable identity, receives a marker at the matched
passage, and moves to the destination project's sidecar; the old sidecar keeps
only a document-local \code{source_moved} tombstone, not a global resolution.
An absent or ambiguous match is retained as \code{needs_anchor} and never
blocks unrelated document rendering.
After a comment deletion is accepted and recorded as a durable tombstone, its unused marker may be removed.

Removing the unique marker of a previously anchored comment from its recorded, readable, managed \code{.tex} source explicitly resolves that comment.
The synchronizer writes a durable global tombstone, removes the record from the live sidecar, prevents older PDF copies from restoring it, and omits it from subsequently published PDFs.
This rule is deliberately guarded: a new or unanchored comment, a missing or unreadable source file, or an identifier that still occurs anywhere in the managed source tree is retained rather than deleted.
A marker moved with its passage remains anchored; duplicate markers remain an ambiguity and are not interpreted as a deletion.
For markup annotations such as highlights, the marker locates the relevant neighborhood while stored selected text and context are used to reconstruct the exact marked range.

The full extraction lifecycle, confidence rules, JSONL fields, and orphan policy are documented in
\TexLink{docs-template/notes_system.tex}{sec:comment-sidecars}{the notes system}.

\section*{Related-work metadata}

\paragraph{\code{\relatedWorkMeta}.}
\agentProposedEdit
  {agent-add-9d1cb594-4f0f-5b3f-aac4-33413e63c489}
  {Add this agent-authored block for review.}
  {}
  {Prints the paper title as a paragraph heading and compact metadata from six arguments: a linked italic paper title, citation information, first-online date, venue, repository or source links, and \code{AuthorsFromZH}.}

\begin{DocCode}
\relatedWorkMeta
  {\href{https://example.com/paper}{\emph{Example Paper}}}
  {23 (queried July 14, 2026)}
  {June 2023}
  {NeurIPS 2023}
  {\url{https://example.com}}
  {-}
\end{DocCode}

A source that defines its own incompatible \code{\relatedWorkMeta} is detected by the renderer and keeps its local definition.

\section*{Bibliographies}

\agentProposedEdit
  {agent-add-1e597d8e-fcc8-581e-9060-b49556f2c5aa}
  {Add this agent-authored block for review.}
  {}
  {The current self-contained policy loads the project-owned database at \path{tex/support/references.bib} or the optional machine-produced override at \path{tex/generated/references.bib}.
Active compilation never obtains a style or bibliography from the workspace-wide \path{tools/latex/} directory.}
\paragraph{\code{\manualBibliography}.}
Adds the References section using the standard \code{plainnat} style.
\agentProposedEdit{agent-remove-bb7d8fc8-732f-4059-98b8-160571c9deb9}{the bibliography database is project-owned rather than loaded from a workspace-wide tools path}{In the normal workspace, it loads the canonical shared database at
\path{tools/latex/tex/support/references.bib}.}{}
\agentProposedEdit{agent-remove-0936333d-c1b0-4fff-bebd-c0a374b0e0ee}{the optional generated override now lives under the canonical tex tree}{If the project contains \path{generated/tex/support/references.bib}, it uses that derived,
self-contained snapshot instead.}{}
\agentProposedEdit
  {agent-add-ab815d3d-1373-5699-9f0d-466f939a7218}
  {Add this agent-authored block for review.}
  {}
  {The command loads \path{tex/support/references.bib} or the optional machine-produced override at \path{tex/generated/references.bib}.}

\begin{DocCode}
% Immediately before \end{document}:
\manualBibliography
\end{DocCode}

\agentProposedEdit{agent-remove-2dc7f1c2-a6e0-4964-85da-4842a61d671d}{generated bibliography inputs are project-owned exact build inputs rather than snapshots of a workspace-wide database}{The generated snapshot is intended for exports such as Overleaf and must be refreshed from the shared database rather than maintained independently.}{}
\agentProposedEdit
  {agent-add-e1732f1e-3098-5ec4-99df-84382b52f15a}
  {Add this agent-authored block for review.}
  {}
  {A generated bibliography override is refreshed by the project's generator and is captured as an exact managed input whenever the target compiles it.}

\section*{Table column type}

\paragraph{\code{P}.}
Defines a fixed-width, left-aligned paragraph column for \code{tabular} and related environments.

\begin{DocCode}
\begin{tabular}{P{0.35\textwidth}P{0.55\textwidth}}
Left column & Right column \\
\end{tabular}
\end{DocCode}

\end{document}
